\scanpage{97}
\mspage{16}
$x_0$ de $\mathcal{T}$, mais sans supposer maintenant $\widetilde{\mathcal{T}}$ connexe. (En pensant au cas où $\widetilde{\mathcal{T}}$ proviendrait d'une situation géométrique $X'\unsure{-|\to}X$, avec $X$ connexe, \deleted{\unsure{et un revêtement}} mais $X'$ seulement $K$-connexe i.e. connexe en tant que schéma sur $K$ --- ce qui signifie qu'il \deleted{\unsure{se}} se décrit en termes galoisiens par $E_0=X'_{x_0}$, sur lequel le produit $\unsure{\frac12|s.}$ direct $\Gamma.\pi_0$ ($\pi_0=\pi_1(X,x_0)=\pi_1(\mathcal{T},x_0)$) opère transitivement, mais sans que l'opération de $\pi_0$ soit forcément transitive.) Associons à $\forall\ x'\in E_0$ son stabilisateur dans $\pi_0$, on trouve une application
\[
(50)\qquad
\diagram{\begin{tikzcd}[ampersand replacement=\&,column sep=large,row sep=small]
E_0 \arrow[r] \& \text{Ens des ss-groupes de }\pi_0 \\
x' \arrow[r,mapsto] \& \text{stabilisateur }\pi_{\unsure{x'|0x'}}
\end{tikzcd}}
\]
qui commute aux actions de $\pi_0$ (op. sur le $2^e$ ens par autom. intérieurs), d'où par passage aux quotients
\[
(51)\qquad
\diagram{\begin{tikzcd}[ampersand replacement=\&,column sep=large]
E_0/\pi_0\overset{\mathrm{def}}=\unsure{\{\pi_0\}|\pi_0}(E_0) \arrow[r,"\overline{\Sigma}:"] \& \operatorname{Ssgr}(\pi)/\pi
\end{tikzcd}}
\]
\inserted{c'est le $\pi_0$ dans le topos des $\pi$-ens.}
D'autre part, $\Theta$ opère sur le $2^{\text{ième}}$ mb de (51) par transport de structure, via \unsure{l'autom.} donné de $\pi$, et \unsure{là} une condition nécessai-
\mspage{96}
